\documentclass{notaaula}

        \titulonota{Ondas e relação fundamental}
        \creditos{Turma 2A · Física · Outubro/2026 · Prof. Josué Cavalcante}

        \begin{document}
        \cabecalho

        \begin{sobre}
        Ondulatória: conceito, natureza das ondas e relação entre rapidez, comprimento de onda e frequência. Habilidades em foco: EM13CNT202, EM13CNT306 e EM13CNT310.
        \end{sobre}

        \section{O que uma onda transporta}

\begin{copiar}{Perturbação que se propaga}
Uma onda é uma perturbação que transporta energia e informação sem transportar matéria de forma
líquida junto com ela. As partículas do meio oscilam em torno de posições de equilíbrio.
\begin{itens}
\item Ondas \dest{mecânicas} precisam de meio: som, onda em corda e onda sísmica.
\item Ondas \dest{eletromagnéticas} propagam-se também no vácuo: rádio, luz e raios X.
\end{itens}
\end{copiar}

\begin{copiar}{Direção da oscilação}
Em onda transversal, a oscilação é perpendicular à propagação. Em onda longitudinal, é paralela.
Uma corda esticada ilustra onda transversal; o som no ar é predominantemente longitudinal.
\atencao{Transversal ou longitudinal descreve a direção da oscilação, não o desenho da trajetória.}
\end{copiar}

\section{Grandezas periódicas}

\begin{copiar}{Amplitude, período e frequência}
A amplitude $A$ mede o afastamento máximo em relação ao equilíbrio. O período $T$ é o tempo de um
ciclo e a frequência $f$ é o número de ciclos por segundo:
\[
  f=\frac{1}{T}.
\]
A unidade de frequência é o hertz: $1\un{Hz}=1\un{s^{-1}}$. O comprimento de onda $\lambda$ é a
distância entre dois pontos consecutivos que oscilam em fase.
\simbolos{$A$ e $\lambda$ (m); $T$ (s); $f$ (Hz).}
\end{copiar}

\begin{exemplo}
Uma fonte realiza 12 oscilações em $3\un{s}$. Então
\[
  f=\frac{12}{3}=4\un{Hz},
  \qquad T=\frac{1}{4}=\resultado{\dec{0,25}\un{s}}.
\]
\end{exemplo}

\section{Relação fundamental}

\begin{copiar}{Rapidez da onda}
Em um período, uma crista avança um comprimento de onda. Portanto,
\[
  v=\frac{\lambda}{T}=\lambda f.
\]
A rapidez é determinada pelas propriedades do meio. Se a fonte muda a frequência no mesmo meio,
o comprimento de onda se ajusta.
\diaadia{Em telecomunicações, diferentes frequências permitem separar canais que se propagam no mesmo espaço.}
\end{copiar}

\begin{exemplo}[Onda em uma corda]
Uma onda tem $\lambda=\dec{0,80}\un{m}$ e $f=5\un{Hz}$.
\[
  v=\lambda f=\dec{0,80}\cdot5
  =\resultado{4\un{m/s}}.
\]
Em $2\un{s}$, a perturbação percorre $8\un{m}$.
\end{exemplo}

\section{Leitura de gráficos}

\begin{copiar}{Espaço e tempo}
Em um retrato $y\times x$, a distância horizontal entre cristas fornece $\lambda$. Em um registro
$y\times t$ feito num ponto, o intervalo entre cristas fornece $T$. O eixo horizontal decide qual
grandeza pode ser lida.
\end{copiar}

\begin{dica}
Amplitude não é comprimento de onda. A amplitude é vertical em um desenho senoidal; $\lambda$ é
a repetição horizontal num retrato espacial.
\end{dica}

\begin{exercicios}
\nivel{1}{Conceitos}
\begin{questoes}
\item O que uma onda transporta?
\item Classifique o som e a luz quanto à necessidade de meio.
\item Diferencie onda transversal de longitudinal.
\item Uma oscilação dura $\dec{0,20}\un{s}$. Qual é a frequência?
\item Em um gráfico $y\times x$, como se identifica $\lambda$?
\nivel{2}{Aplicação}
\item Uma onda tem $f=10\un{Hz}$ e $\lambda=\dec{0,50}\un{m}$. Calcule $v$.
\item Uma onda se propaga a $12\un{m/s}$ com $f=3\un{Hz}$. Ache $\lambda$.
\item Uma fonte faz 90 ciclos em $30\un{s}$. Ache $f$ e $T$.
\item No mesmo meio, a frequência dobra. O que acontece com $\lambda$?
\nivel{3}{Síntese}
\item Uma onda percorre $150\un{m}$ em $\dec{0,50}\un{s}$ e tem $f=100\un{Hz}$. Ache $v$ e $\lambda$.
\item Explique por que as partículas de uma corda não viajam junto com a onda até a outra extremidade.
\item Compare o que se lê em gráficos $y\times x$ e $y\times t$.
\end{questoes}
\gabarito{4) $5\un{Hz}$; 6) $5\un{m/s}$; 7) $4\un{m}$;
8) $3\un{Hz}$ e $1/3\un{s}$; 9) cai à metade; 10) $300\un{m/s}$ e $3\un{m}$.}
\end{exercicios}


\end{document}
